% Relatorio do Trabalho Pratico -- Computacao Distribuida (PUC Minas, 2026/2)
%
% ESTE ARQUIVO E UM ANDAIME: a estrutura, as figuras e a tabela estao prontas;
% o texto de cada secao e escrito pelo grupo (o enunciado desconsidera texto
% gerado automaticamente). Os marcadores \pendente{...} mostram o que falta.
%
% Limite da Entrega 1: 4 paginas. Limite da Entrega 2: 6 paginas.
% Template: modelo de artigo da SBC (sbc-template.sty vendorizado nesta pasta).
%
% Compilar (nao precisa de TeX no host):
%   docker run --rm -v "$PWD":/w -w /w texlive/texlive:latest \
%     latexmk -pdf -interaction=nonstopmode -halt-on-error main.tex

\documentclass[12pt]{article}

\usepackage{sbc-template}
\usepackage{graphicx}
\usepackage{url}
\usepackage[utf8]{inputenc}
\usepackage[T1]{fontenc}
\usepackage[brazil]{babel}
\usepackage{booktabs}
\usepackage{xcolor}
\usepackage{tikz}
\usetikzlibrary{arrows.meta,positioning,fit,calc}
\usepackage{needspace}

% Marcador visivel de texto que ainda precisa ser escrito pelo grupo.
% Troque para \pendentesfalse para ver o PDF sem os marcadores e medir o espaco livre
% (foi assim que conferimos que ainda sobram ~1,5 pagina para a prosa).
\newif\ifpendentes \pendentestrue
\newcommand{\pendente}[1]{\ifpendentes\textcolor{red}{\textbf{[PENDENTE: #1]}}\fi}

\sloppy

\title{Col\^onia de Formigas Distribu\'ida para o Caixeiro-Viajante com Elei\c{c}\~ao de L\'ider, Rel\'ogio de Lamport e Exclus\~ao M\'utua}

% R1: autores confirmados (o enunciado pede grupo de 4 a 5).
\author{Pedro Henrique Vila\c{c}a Valverde\inst{1}, Marcus Prado Silva\inst{1},\\
Gabriel Batista de Almeida\inst{1}, Eduardo Alves Salgado Lisboa Oliveira\inst{1}}

\address{Instituto de Ci\^encias Exatas e Inform\'atica (ICEx) -- PUC Minas\\
Belo Horizonte -- MG -- Brasil
}

\begin{document}

\maketitle

\begin{resumo}
% Regra do template: no maximo 10 linhas, na primeira pagina.
% ROTEIRO (4 frases, escrever por ultimo):
%  1) problema e abordagem: "Este trabalho apresenta uma colonia de formigas distribuida
%     para o caixeiro-viajante, executada em nos Java que trocam mensagens por gRPC."
%  2) como os nos se coordenam: lider eleito por Bully, relogio de Lamport, exclusao mutua
%     de Ricart-Agrawala para a matriz de feromonio.
%  3) o que foi medido: 40 testes, convergencia das replicas, failover em ~7,5 s.
%  4) conclusao em 1 frase: o que funcionou e o que fica para a proxima etapa.
\pendente{resumo em portugu\^es, at\'e 10 linhas}
\end{resumo}

\begin{abstract}
% Traducao do resumo, tambem ate 10 linhas (regra do template).
\pendente{abstract em ingl\^es, at\'e 10 linhas}
\end{abstract}

\noindent\textbf{Palavras-chave:} \pendente{3 a 5: computa\c{c}\~ao distribu\'ida; col\^onia de formigas;
elei\c{c}\~ao de l\'ider; exclus\~ao m\'utua}

\section{Introdu\c{c}\~ao}

% FONTES (nao copiar: usar como base factual)
%   docs/algoritmos/aco-tsp.md ....... problema, ACO, instancias
%   docs/ENUNCIADO.md ................ requisitos e criterios de avaliacao
%   README.md ........................ como o sistema roda

\pendente{problema: TSP e por que distribuir}

\pendente{o que foi construido, em 3 linhas}

\pendente{4 requisitos atendidos (o enunciado exige 2) -- dizer explicitamente}

\pendente{organizacao do artigo}

\section{Arquitetura do Sistema}

% FONTES: docs/ARQUITETURA.md (secoes 1 a 4), docs/rede-grpc.md, docs/rodadas.md,
%         docs/concorrencia.md, docs/LAN.md (implantacao)

\pendente{vis\~ao geral: N n\'os sim\'etricos, mesmo bin\'ario, papel definido pela elei\c{c}\~ao}

\begin{figure}[ht]
  \centering
  \scalebox{0.88}{% Figura 1 -- entidades e principais mensagens (TikZ, vetorial; nao depende de PNG externo).
% Uso: % Figura 1 -- entidades e principais mensagens (TikZ, vetorial; nao depende de PNG externo).
% Uso: \input{figuras/arquitetura} dentro de um ambiente figure.
\begin{tikzpicture}[
    font=\footnotesize,
    no/.style={draw, rounded corners=2pt, minimum width=1.45cm, minimum height=0.75cm, align=center},
    lider/.style={no, fill=black!12, very thick},
    seta/.style={-{Latex[length=1.7mm]}, semithick},
    rot/.style={font=\scriptsize, inner sep=1.5pt}
  ]

  % Lider
  \node[lider] (n4) at (1.55,2.15) {N4\\[-1pt]\scriptsize l\'ider};

  % Trabalhadores
  \node[no] (n1) at (0,0)   {N1};
  \node[no] (n2) at (1.55,0) {N2};
  \node[no] (n3) at (3.10,0) {N3};

  % Lider -> nos (StartRound) e nos -> lider (RoundDone / Heartbeat)
  \foreach \n in {n1,n2,n3} {
      \draw[seta] (n4) to[bend left=10] (\n);
      \draw[seta] (\n) to[bend left=10] (n4);
    }

  % Legenda (evita rotulos em cima das setas)
  \draw[seta] (0.0,-1.20) -- ++(0.55,0) node[right,rot] {StartRound: l\'ider $\rightarrow$ n\'os};
  \draw[seta] (0.0,-1.60) -- ++(0.55,0) node[right,rot] {RoundDone e Heartbeat: n\'o $\rightarrow$ l\'ider};
  \draw[seta, dashed] (0.0,-2.00) -- ++(0.55,0) node[right,rot]
    {UpdatePheromone: n\'o na se\c{c}\~ao cr\'itica $\rightarrow$ todos (matriz $n\times n$)};

\end{tikzpicture}
 dentro de um ambiente figure.
\begin{tikzpicture}[
    font=\footnotesize,
    no/.style={draw, rounded corners=2pt, minimum width=1.45cm, minimum height=0.75cm, align=center},
    lider/.style={no, fill=black!12, very thick},
    seta/.style={-{Latex[length=1.7mm]}, semithick},
    rot/.style={font=\scriptsize, inner sep=1.5pt}
  ]

  % Lider
  \node[lider] (n4) at (1.55,2.15) {N4\\[-1pt]\scriptsize l\'ider};

  % Trabalhadores
  \node[no] (n1) at (0,0)   {N1};
  \node[no] (n2) at (1.55,0) {N2};
  \node[no] (n3) at (3.10,0) {N3};

  % Lider -> nos (StartRound) e nos -> lider (RoundDone / Heartbeat)
  \foreach \n in {n1,n2,n3} {
      \draw[seta] (n4) to[bend left=10] (\n);
      \draw[seta] (\n) to[bend left=10] (n4);
    }

  % Legenda (evita rotulos em cima das setas)
  \draw[seta] (0.0,-1.20) -- ++(0.55,0) node[right,rot] {StartRound: l\'ider $\rightarrow$ n\'os};
  \draw[seta] (0.0,-1.60) -- ++(0.55,0) node[right,rot] {RoundDone e Heartbeat: n\'o $\rightarrow$ l\'ider};
  \draw[seta, dashed] (0.0,-2.00) -- ++(0.55,0) node[right,rot]
    {UpdatePheromone: n\'o na se\c{c}\~ao cr\'itica $\rightarrow$ todos (matriz $n\times n$)};

\end{tikzpicture}
}
  \caption{Entidades e mensagens. Todos os n\'os s\~ao trabalhadores; o l\'ider (maior id vivo)
  coordena as rodadas. A matriz de ferom\^onio \'e replicada e s\'o muda dentro da se\c{c}\~ao
  cr\'itica.}
  \label{fig:arquitetura}
\end{figure}

\subsection{Entidades e pap\'eis}

\pendente{l\'ider, trabalhador, detector de l\'ider e m\'odulos de suporte}

\subsection{Mensagens}

% Tabela factual extraida de src/main/proto/node.proto e de rpc/PeerNetwork.
\begin{table}[ht]
  \centering
  \footnotesize
  \begin{tabular}{@{}lll@{}}
    \toprule
    Mensagem & Dire\c{c}\~ao & Papel \\
    \midrule
    \texttt{Heartbeat}            & n\'o $\to$ l\'ider       & detector de falha e handshake de configura\c{c}\~ao \\
    \texttt{Election} / \texttt{Coordinator} & todos $\to$ todos & elei\c{c}\~ao (Bully) \\
    \texttt{RequestCS} / \texttt{ReplyCS}    & n\'o $\leftrightarrow$ n\'o & Ricart-Agrawala (com \texttt{request\_ts}) \\
    \texttt{StartRound}           & l\'ider $\to$ todos      & in\'icio de rodada \\
    \texttt{RoundDone}            & n\'o $\to$ coordenador   & melhor rota local da rodada \\
    \texttt{UpdatePheromone}      & n\'o $\to$ todos         & matriz $n\times n$ + melhor rota (s\'incrono) \\
    \bottomrule
  \end{tabular}
  \caption{Mensagens do servi\c{c}o \texttt{NodeService} (gRPC). Toda mensagem carrega
  \texttt{sender\_id} e um carimbo de Lamport no cabe\c{c}alho, al\'em do token do cluster.}
  \label{tab:mensagens}
\end{table}

\pendente{por que um servi\c{c}o \'unico e o que vai no cabe\c{c}alho}

\subsection{Uma rodada}

\pendente{formigas $\to$ se\c{c}\~ao cr\'itica (atualiza + propaga) $\to$ RoundDone; o que roda fora da SC}

\section{Requisitos Implementados}

% FONTES: docs/ARQUITETURA.md §5 (5.1 gRPC, 5.2 Bully, 5.3 Lamport, 5.4 Ricart-Agrawala),
%         docs/algoritmos/*.md (algoritmo + diagrama de sequencia de cada um),
%         docs/classes/*.md (detalhe de implementacao por classe)
% O enunciado pede NO MINIMO DOIS requisitos; aqui sao quatro. Cada subsecao deve caber em
% ~1/3 de pagina: algoritmo, mensagens e UMA decisao de projeto.

\subsection{RPC -- gRPC}
% src/main/proto/node.proto, rpc/PeerNetwork.java, rpc/NodeServiceImpl.java, rpc/ClusterAuth.java
\pendente{servi\c{c}o \'unico; deadlines; token no interceptor; 64 MiB}

\subsection{Elei\c{c}\~ao de l\'ider -- Bully}
% election/BullyElection.java, election/LeaderMonitor.java, docs/algoritmos/bully.md
\pendente{3 passos; maior id; 5 heartbeats; contesta\c{c}\~ao de id menor; failover (Tabela 2)}

\subsection{Rel\'ogio l\'ogico -- Lamport}
% clock/LamportClock.java, docs/algoritmos/lamport.md
\pendente{tick no envio, max+1 no recebimento; ordena pedidos e numera logs}

\subsection{Exclus\~ao m\'utua -- Ricart-Agrawala}
% mutex/RicartAgrawala.java, docs/algoritmos/ricart-agrawala.md
\pendente{RELEASED/WANTED/HELD; prioridade (ts, id); REPLY carrega request\_ts; 2(N-1) mensagens;
falha = permiss\~ao (risco assumido, vira desafio na se\c{c}\~ao 5)}

\begin{figure}[t]
  \centering
  \scalebox{0.85}{% Figura 2 -- Ricart-Agrawala com tres participantes.
% Uso: % Figura 2 -- Ricart-Agrawala com tres participantes.
% Uso: \input{figuras/sequencia-ra} dentro de um ambiente figure.
\begin{tikzpicture}[
    font=\footnotesize,
    cab/.style={draw, rounded corners=2pt, minimum width=1.15cm, minimum height=0.6cm, fill=black!6},
    seta/.style={-{Latex[length=1.7mm]}, semithick},
    rot/.style={font=\scriptsize, inner sep=1.5pt, fill=white}
  ]

  % Cabecalhos e linhas de vida
  \node[cab] (h1) at (0,0)   {N1};
  \node[cab] (h2) at (3.2,0) {N2};
  \node[cab] (h3) at (6.4,0) {N3};
  \foreach \x in {0,3.2,6.4} { \draw[dashed, gray] (\x,-0.32) -- (\x,-4.45); }

  % REQUEST do N1 para os dois peers
  \draw[seta] (0,-1.00) -- node[rot, above] {REQUEST $(ts=5)$} (3.2,-1.00);
  \draw[seta] (0,-1.55) -- node[rot, above] {REQUEST $(ts=5)$} (6.4,-1.55);

  % Permissoes
  \draw[seta] (3.2,-2.10) -- node[rot, above] {REPLY} (0,-2.10);
  \draw[seta] (6.4,-2.55) -- node[rot, above] {REPLY} (0,-2.55);

  % N1 entra na secao critica com as duas permissoes
  \fill[black!12] (-0.42,-3.00) rectangle (0.42,-4.15);
  \node[font=\scriptsize, align=center, text width=1.05cm] at (0,-3.58)
    {se\c{c}\~ao\\cr\'itica};

\end{tikzpicture}
 dentro de um ambiente figure.
\begin{tikzpicture}[
    font=\footnotesize,
    cab/.style={draw, rounded corners=2pt, minimum width=1.15cm, minimum height=0.6cm, fill=black!6},
    seta/.style={-{Latex[length=1.7mm]}, semithick},
    rot/.style={font=\scriptsize, inner sep=1.5pt, fill=white}
  ]

  % Cabecalhos e linhas de vida
  \node[cab] (h1) at (0,0)   {N1};
  \node[cab] (h2) at (3.2,0) {N2};
  \node[cab] (h3) at (6.4,0) {N3};
  \foreach \x in {0,3.2,6.4} { \draw[dashed, gray] (\x,-0.32) -- (\x,-4.45); }

  % REQUEST do N1 para os dois peers
  \draw[seta] (0,-1.00) -- node[rot, above] {REQUEST $(ts=5)$} (3.2,-1.00);
  \draw[seta] (0,-1.55) -- node[rot, above] {REQUEST $(ts=5)$} (6.4,-1.55);

  % Permissoes
  \draw[seta] (3.2,-2.10) -- node[rot, above] {REPLY} (0,-2.10);
  \draw[seta] (6.4,-2.55) -- node[rot, above] {REPLY} (0,-2.55);

  % N1 entra na secao critica com as duas permissoes
  \fill[black!12] (-0.42,-3.00) rectangle (0.42,-4.15);
  \node[font=\scriptsize, align=center, text width=1.05cm] at (0,-3.58)
    {se\c{c}\~ao\\cr\'itica};

\end{tikzpicture}
}
  \caption{Ricart-Agrawala com tr\^es participantes: N1 s\'o entra na se\c{c}\~ao cr\'itica depois
  de receber as duas permiss\~oes. O par $(ts, id)$ decide quem espera: um pedido mais novo pode
  ser adiado e respondido s\'o na sa\'ida da se\c{c}\~ao cr\'itica.}
  \label{fig:ra}
\end{figure}

\section{Resultados}

% ============================ ROTEIRO ============================
% Fontes dos numeros: docs/AUDITORIA.md (20/10) e ensaio de compose LAN em 1 maquina (28/10).
% R4: confirmar com o ensaio de 2 maquinas (docs/evidencias/) ou remover a frase sobre
% maquinas fisicas.
%
% [4.1 O que a tabela mostra -- 1 paragrafo]
%  - Replicas convergem: ao fim das rodadas as matrizes sao identicas nos 4 nos
%    (ClusterIntegrationTest.elegeMaiorIdExecutaTodasAsRodadasEReplicasConvergem).
%  - Custo por rodada: ~50 ms de computo local contra 1 broadcast sincrono; ou seja, o
%    gargalo e a rede, nao as formigas.
%  - Failover: ~7,5 s na LAN = 5 heartbeats de 1,5 s + eleicao; no compose (heartbeat 1 s)
%    sao ~5 s. Comparar com o que se esperaria sem detector (rodada de 30 s por timeout).
%  - Honestidade metodologica: a linha de tempo foi medida em 1 maquina com 4 containers;
%    4000 rodadas e 600 rodadas com seguidor morto vem da auditoria de 20/10.
%
% [4.2 Ensaio em 2 maquinas -- 1 paragrafo, so se for feito]
%  - Mesmo commit nas duas, porta 50051 liberada, mesmo CLUSTER_TOKEN.
%  - Reportar: 30 rodadas, ms/rodada, melhor rota igual nos dois logs, tempo de failover.
%  - Se nao for feito, APAGAR a frase e nao mencionar maquinas fisicas.
% =================================================================

\begin{table}[t]
  \centering
  \footnotesize
  \begin{tabular}{@{}ll@{}}
    \toprule
    Medida & Valor observado \\
    \midrule
    Testes automatizados (unit\'arios + integra\c{c}\~ao) & 40, todos passando \\
    Rodadas de 30 (60 cidades, 40 formigas por n\'o)      & $\approx$1,5 s no total ($\approx$50 ms/rodada) \\
    Melhor rota encontrada (seed 42)                       & 6067 \\
    Failover do l\'ider (4 n\'os, configura\c{c}\~ao LAN)  & $\approx$7,5 s (5 heartbeats de 1,5 s + elei\c{c}\~ao) \\
    Estabilidade em execu\c{c}\~ao longa                   & 4000 rodadas sem timeout \\
    Queda de seguidor no meio da execu\c{c}\~ao            & sobreviventes conclu\'iram as 600 rodadas \\
    \bottomrule
  \end{tabular}
  \caption{Medidas do sistema. A linha de tempo por rodada foi obtida em uma \'unica m\'aquina
  (4 cont\^eineres); os n\'umeros de execu\c{c}\~oes longas v\^em da auditoria de 20/10.}
  \label{tab:resultados}
\end{table}

\pendente{o que a tabela mostra: r\'eplicas convergem, custo de rede por rodada, tempo de detec\c{c}\~ao}

\pendente{ensaio em 2 m\'aquinas f\'isicas (se j\'a feito; sen\~ao remover esta frase)}

\section{Desafios e Problemas Encontrados}

O detector trabalha por \textit{timeout} e não distingue nó lento de nó morto. Quando um
\texttt{REQUEST} falha ou expira, o nó trata o silêncio como permissão: é o que impede o
\textit{cluster} de travar quando um \textit{peer} morre, mas o mesmo caminho dispara quando ele
apenas sofre uma pausa de GC. Reproduzimos o caso com dois nós reais, e ambos entraram na seção
crítica. Não corrigimos: separar lento de morto exige quórum ou consenso, e privilegiamos a
disponibilidade.

Já a queda de um \textit{peer} entre o \texttt{REQUEST} e o \texttt{REPLY} deixava a aquisição
presa, pois o nó aguarda $N-1$ permissões; passamos a reenviar o pedido a cada 500 ms e a
recontatar quem volta após a quarentena, o que também destravou a queda de seguidor, que custava
30 s por rodada. Configurações divergentes, por fim, geravam \textit{split-brain} silencioso e
agora são recusadas na subida, com código 2.

\section{Papel de Cada Aluno}

A divisão seguiu os quatro pacotes do sistema. Pedro Henrique Vilaça Valverde ficou com a camada
de comunicação e infraestrutura: prazos de rede local, \textit{handshake} de token, leitor de
TSPLIB e contêineres. Marcus Prado Silva assumiu o Ricart-Agrawala, seus testes de integração e a
documentação técnica. Gabriel Batista de Almeida respondeu pelo ACO e pelas rodadas: formigas,
matriz de feromônio e ciclo de rodadas. Eduardo Alves Salgado Lisboa Oliveira respondeu pela
eleição: o algoritmo Bully, o detector de \textit{heartbeat} e o roteiro de \textit{failover}.


% Bibliografia manual (6 entradas) para compilar sem precisar do sbc.bst oficial.
% Para usar BibTeX: baixe sbc.bst, troque este bloco por
%   \bibliographystyle{sbc} \bibliography{referencias}
\Needspace*{12\baselineskip} % nao deixa a lista de referencias quebrar entre paginas
% Bibliografia no formato da SBC, escrita a mao para dispensar o sbc.bst na compilacao.
% As mesmas entradas estao em referencias.bib, prontas para o BibTeX se o grupo preferir
% (ver README.md desta pasta).

\begin{thebibliography}{5}

\bibitem{lamport78}
LAMPORT, L. Time, clocks, and the ordering of events in a distributed system.
\emph{Communications of the ACM}, v.~21, n.~7, p.~558--565, 1978.

\bibitem{ricart81}
RICART, G.; AGRAWALA, A. K. An optimal algorithm for mutual exclusion in computer networks.
\emph{Communications of the ACM}, v.~24, n.~1, p.~9--17, 1981.

\bibitem{garcia82}
GARCIA-MOLINA, H. Elections in a distributed computing system.
\emph{IEEE Transactions on Computers}, v.~C-31, n.~1, p.~48--59, 1982.

\bibitem{dorigo97}
DORIGO, M.; GAMBARDELLA, L. M. Ant colony system: a cooperative learning approach to the
traveling salesman problem. \emph{IEEE Transactions on Evolutionary Computation}, v.~1, n.~1,
p.~53--66, 1997.

\bibitem{reinelt91}
REINELT, G. TSPLIB -- a traveling salesman problem library.
\emph{ORSA Journal on Computing}, v.~3, n.~4, p.~376--384, 1991.

\end{thebibliography}


\end{document}
